\textbf{Phần 1. Trắc nghiệm lựa chọn}
\setcounterpageref{ex}{0}

\begin{ex} %[1D10V41]
Tính giá trị của biểu thức: $42 + 35 - 24$.
\choice
{$52$}
{$54$}
{\True $53$}
{$51$}
\loigiai{
Thực hiện: $42 + 35 = 77$; $77 - 24 = 53$.
\textbf{Đáp số:} $53$.\par
\textbf{Đáp án của câu hỏi là C}
}
\end{ex}

\begin{ex} %[1D10V41]
Tính giá trị của biểu thức: $88 - 53 + 14$.
\choice
{$48$}
{\True $49$}
{$50$}
{$47$}
\loigiai{
Thực hiện: $88 - 53 = 35$; $35 + 14 = 49$.
\textbf{Đáp số:} $49$.\par
\textbf{Đáp án của câu hỏi là B}
}
\end{ex}

\begin{ex} %[1D10V41]
Tính giá trị của biểu thức: $99 - 44 - 33$.
\choice
{$21$}
{$23$}
{$20$}
{\True $22$}
\loigiai{
Thực hiện: $99 - 44 = 55$; $55 - 33 = 22$.
\textbf{Đáp số:} $22$.\par
\textbf{Đáp án của câu hỏi là D}
}
\end{ex}

\begin{ex} %[1D10V41]
Tìm số tự nhiên lớn nhất có hai chữ số.
\choice
{\True $99$}
{$98$}
{$100$}
{$97$}
\loigiai{
Số tự nhiên lớn nhất có hai chữ số là số $99$.
\textbf{Đáp số:} $99$.\par
\textbf{Đáp án của câu hỏi là A}
}
\end{ex}

\begin{ex} %[1D10V41]
Tìm số tự nhiên bé nhất có hai chữ số khác nhau.
\choice
{$0$}
{$20$}
{\True $10$}
{$9$}
\loigiai{
Số bé nhất có hai chữ số khác nhau là số $10$.
\textbf{Đáp số:} $10$.\par
\textbf{Đáp án của câu hỏi là C}
}
\end{ex}

\begin{ex} %[1D10V41]
Tìm số liền sau của số 99.
\choice
{$90$}
{\True $100$}
{$110$}
{$99$}
\loigiai{
Số liền sau của 99 là số $100$.
\textbf{Đáp số:} $100$.\par
\textbf{Đáp án của câu hỏi là B}
}
\end{ex}

\begin{ex} %[1D10V41]
Tìm số liền trước của số 80.
\choice
{$78$}
{$80$}
{$77$}
{\True $79$}
\loigiai{
Số liền trước của 80 là số $79$.
\textbf{Đáp số:} $79$.\par
\textbf{Đáp án của câu hỏi là D}
}
\end{ex}

\begin{ex} %[1D10V41]
Điền số thích hợp: $54 + \dots = 89$.
\choice
{\True $35$}
{$34$}
{$36$}
{$33$}
\loigiai{
Số cần điền là: $89 - 54 = 35$.
\textbf{Đáp số:} $35$.\par
\textbf{Đáp án của câu hỏi là A}
}
\end{ex}

\begin{ex} %[1D10V41]
Điền số thích hợp: $76 - \dots = 42$.
\choice
{$33$}
{$35$}
{\True $34$}
{$32$}
\loigiai{
Số cần điền là: $76 - 42 = 34$.
\textbf{Đáp số:} $34$.\par
\textbf{Đáp án của câu hỏi là C}
}
\end{ex}

\begin{ex} %[1D10V41]
Một trường có 55 học sinh nam và 43 học sinh nữ khối 1. Hỏi khối 1 có tất cả bao nhiêu học sinh?
\choice
{$97$}
{\True $98$}
{$99$}
{$96$}
\loigiai{
Tổng số học sinh khối 1 là: $55 + 43 = 98$ (học sinh).
\textbf{Đáp số:} $98$.\par
\textbf{Đáp án của câu hỏi là B}
}
\end{ex}

\begin{ex} %[1D10V41]
Số học sinh nam nhiều hơn số học sinh nữ ở bài toán trên là bao nhiêu bạn?
\choice
{$11$}
{$13$}
{$10$}
{\True $12$}
\loigiai{
Số học sinh nam nhiều hơn nữ là: $55 - 43 = 12$ (học sinh).
\textbf{Đáp số:} $12$.\par
\textbf{Đáp án của câu hỏi là D}
}
\end{ex}

\begin{ex} %[1D10V41]
Một cửa hàng có 85 chiếc xe đạp, đã bán được 53 chiếc xe đạp. Hỏi cửa hàng còn lại bao nhiêu chiếc xe đạp?
\choice
{\True $32$}
{$31$}
{$33$}
{$30$}
\loigiai{
Số xe đạp còn lại là: $85 - 53 = 32$ (chiếc xe đạp).
\textbf{Đáp số:} $32$.\par
\textbf{Đáp án của câu hỏi là A}
}
\end{ex}

\textbf{Phần 2. Trắc nghiệm đúng sai}
\setcounterpageref{ex}{0}

\begin{ex} %[1D10V41]
Xét tính đúng hoặc sai của các khẳng định toán học sau đây:
\choiceTF
{\True Kết quả của bài toán «Một sợi dây dài 95 cm. Người ta cắt đi hai đoạn dài 30 cm và 25 cm. Sợi dây còn lại dài bao nhiêu xăng-ti-mét» là $40\text{ cm}$.}
{Kết quả của bài toán «Đồng hồ chỉ 10 giờ đêm (tức 22 giờ). Sau 2 giờ nữa là mấy giờ đêm (ngày mới)» là $13\text{ giờ}$.}
{Kết quả của bài toán «Nếu hôm nay là Thứ Năm ngày 14 thì Thứ Năm tuần trước là ngày mùng mấy» là $6\text{ ngày}$.}
{\True Kết quả của bài toán «Nếu hôm nay là Thứ Bảy ngày 21 thì Thứ Bảy tuần sau là ngày bao nhiêu» là $28\text{ ngày}$.}
\loigiai{
\begin{itemize}
    \item Ý a) Đúng: Tổng độ dài đã cắt: $30 + 25 = 55$ (cm).\nĐộ dài còn lại là: $95 - 55 = 40$ (cm).
\textbf{Đáp số:} $40$.
    \item Ý b) Sai: Kết quả đúng là $12$, không phải $13$. Lúc đó là: $10 + 2 = 12$ (giờ đêm, tức 0 giờ ngày mới).
\textbf{Đáp số:} $12$.
    \item Ý c) Sai: Kết quả đúng là $7$, không phải $6$. Thứ Năm tuần trước là ngày: $14 - 7 = 7$.
\textbf{Đáp số:} $7$.
    \item Ý d) Đúng: Thứ Bảy tuần sau là ngày: $21 + 7 = 28$.
\textbf{Đáp số:} $28$.
\end{itemize}
\begin{itemize}
    \item a) \textbf{Đúng}.
    \item b) \textbf{Sai}.
    \item c) \textbf{Sai}.
    \item d) \textbf{Đúng}.
\end{itemize}
}
\end{ex}

\begin{ex} %[1D10V41]
Xét tính đúng hoặc sai của các bài toán và phép tính sau đây:
\choiceTF
{Kết quả của bài toán «Một hình chữ nhật có chiều dài 30 cm, chiều rộng 20 cm. Chiều dài hơn chiều rộng bao nhiêu xăng-ti-mét» là $11\text{ cm}$.}
{\True Khối lập phương có $6$ mặt phẳng.}
{\True Khối hộp chữ nhật có $12$ cạnh.}
{Kết quả của bài toán «Nếu hôm nay là Thứ Năm ngày 14 thì Thứ Năm tuần trước là ngày mùng mấy» là $6\text{ ngày}$.}
\loigiai{
\begin{itemize}
    \item Ý a) Sai: Kết quả đúng là $10$, không phải $11$. Chiều dài hơn chiều rộng là: $30 - 20 = 10$ (cm).
\textbf{Đáp số:} $10$.
    \item Ý b) Đúng: Khối lập phương có đúng 6 mặt phẳng.
\textbf{Đáp số:} $6$.
    \item Ý c) Đúng: Khối hộp chữ nhật có đúng 12 cạnh.
\textbf{Đáp số:} $12$.
    \item Ý d) Sai: Kết quả đúng là $7$, không phải $6$. Thứ Năm tuần trước là ngày: $14 - 7 = 7$.
\textbf{Đáp số:} $7$.
\end{itemize}
\begin{itemize}
    \item a) \textbf{Sai}.
    \item b) \textbf{Đúng}.
    \item c) \textbf{Đúng}.
    \item d) \textbf{Sai}.
\end{itemize}
}
\end{ex}

\textbf{Phần 3. Trả lời ngắn}
\setcounterpageref{ex}{0}

\begin{ex} %[1D10V41]
Tìm hai số có tổng bằng 77 và hiệu bằng 11. Số lớn hơn là số mấy?
\shortans{44}
\loigiai{
\textbf{Đáp số:} $44$\par
Hai số là 44 và 33 ($44 + 33 = 77$ và $44 - 33 = 11$). Số lớn hơn là $44$.
\textbf{Đáp số:} $44$.
}
\end{ex}

\begin{ex} %[1D10V41]
Tìm hai số có tổng bằng 88 và hiệu bằng 22. Số bé hơn là số mấy?
\shortans{33}
\loigiai{
\textbf{Đáp số:} $33$\par
Hai số là 55 và 33 ($55 + 33 = 88$ và $55 - 33 = 22$). Số bé hơn là $33$.
\textbf{Đáp số:} $33$.
}
\end{ex}

\begin{ex} %[1D10V41]
Có bao nhiêu số tròn chục có hai chữ số?
\shortans{9}
\loigiai{
\textbf{Đáp số:} $9$\par
Có 9 số tròn chục: $10, 20, 30, 40, 50, 60, 70, 80, 90$.
\textbf{Đáp số:} $9$.
}
\end{ex}

\begin{ex} %[1D10V41]
Có bao nhiêu số có hai chữ số giống nhau?
\shortans{9}
\loigiai{
\textbf{Đáp số:} $9$\par
Có 9 số: $11, 22, 33, 44, 55, 66, 77, 88, 99$.
\textbf{Đáp số:} $9$.
}
\end{ex}

\begin{ex} %[1D10V41]
Một cuộn ruy băng dài 90 cm. Bạn Linh cắt lấy 3 đoạn, mỗi đoạn dài 20 cm. Hỏi cuộn ruy băng còn lại dài bao nhiêu xăng-ti-mét?
\shortans{30}
\loigiai{
\textbf{Đáp số:} $30$\par
Tổng độ dài 3 đoạn cắt là: $20 + 20 + 20 = 60$ (cm).\nĐộ dài còn lại là: $90 - 60 = 30$ (cm).
\textbf{Đáp số:} $30$.
}
\end{ex}

\begin{ex} %[1D10V41]
Tính giá trị của biểu thức: $99 - 11 - 22 - 33$.
\shortans{33}
\loigiai{
\textbf{Đáp số:} $33$\par
Thực hiện: $99 - 11 = 88$; $88 - 22 = 66$; $66 - 33 = 33$.
\textbf{Đáp số:} $33$.
}
\end{ex}
